\section{Flatten}

\paragraph{Input and output.}
Flatten consumes an elaborated closed AST and emits a flat closed Axon AST.

\paragraph{Contract.}
The flat program has no statement-level \code{if}, no \code{for}, no
\code{scope}, no relative paths, and no hidden callee path sugar.  Complex
expressions are lowered into explicit bind statements, loop constructs are
converted into helper definitions, and returns/yields are restricted to atomic
variables or tuples of variables.

\paragraph{Path semantics.}
Flatten resolves lexical scopes into explicit absolute or templated absolute
paths.  Scoped callees receive ordinary \code{Path}-typed parameters rather
than syntactic path-sugar parameters.

\paragraph{Rationale.}
Flatten produces the small core language expected by typecheck2 and Graph IR
lowering.  Backends should not implement surface Axon control constructs.

\paragraph{Formal view.}
Flatten maps elaborated closed ASTs to flat closed ASTs:
\[
  \mathsf{flatten}: A_e \rightarrow A_{\mathrm{flat}}.
\]
It preserves evaluation order by introducing fresh bind names.  If
\(\mathsf{atoms}(e)\) are variables, literals, paths, or tuples of atoms, then
flat bind right-hand sides are calls or primitive expressions over atoms.

\paragraph{Example: expression flattening.}
The nested source expression
\begin{verbatim}
y <- NN.linear @w (Tensor.reshape x shape=[B, S, D]) dim=O
\end{verbatim}
becomes a sequence such as
\begin{verbatim}
_v1 <- Tensor.reshape x shape=[B, S, D]
y <- NN.linear @@w _v1 dim=O
\end{verbatim}
where the generated temporary preserves eager evaluation order.

\paragraph{Example: scope and loop lowering.}
A scoped loop
\begin{verbatim}
x <- for@h i <- [0..NUM_LAYERS) carry (x) do
  x <- scope@attn do
    NN.linear@c_proj x
  yield x
\end{verbatim}
is converted into helper definitions and calls with ordinary path parameters.
The resulting paths are absolute templates such as
\code{@@'h.\{i\}.attn.c\_proj'} rather than dynamic scope statements.
